\documentclass[12pt]{scrartcl}


%Matthis Vorlage
\usepackage{../mathe_definition}
%wird in Matthis Volage auch gemacht aber das rafft mein texstudio nicht und highlighted Dinge falsch
\usepackage{amsmath} 

%Ausgelagerter Code fürs schöne seitenlayout
\date{\today}
\newcommand{\Blattnummer}{%
	4%
}

\title{{\itshape Stochastische Prozesse}}
\module{Blatt 1}
\author{Philip Wunderlich \and Jonathan Langer}
%\matriculation{3737927\and 3750030}

\fancypagestyle{netter_header}{
	\renewcommand{\and}{\\}
	\fancyhead[L]{\MyAuthor}
	\fancyhead[C]{\MyTitle\\ \MyModule}
	\fancyhead[R]{\MyDate\\Seite \thepage{} von \pageref{LastPage}}
	\fancyfoot[C]{}
}
\pagestyle{netter_header}

%Abkürzungen
\newcommand{\einruecken}{\begin{addmargin}[25pt]{0pt}} 
	%LaTeX mag keine Umlaute :(
	\newcommand{\ausruecken}{\end{addmargin}}
\newcommand{\R}{\ensuremath{\mathbb{R}}}
\newcommand{\Om}{\ensuremath{\Omega}}
\newcommand{\om}{\ensuremath{\omega}}
\newcommand{\A}{\ensuremath{\mathcal{A}}}
\newcommand{\B}{\ensuremath{\mathcal{B}}}
\newcommand{\C}{\ensuremath{\mathcal{C}}}
\newcommand{\WraumAnn}{\ensuremath{(\Om, \A, P)} ein Wahrscheinlichkeitsraum}
\newcommand{\MraumAnn}{\ensuremath{(E,\B)} ein Messraum}
\newcommand{\Xt}{\ensuremath{(X_t)_{t\in T}}}
\newcommand{\Yt}{\ensuremath{(Y_t)_{t\in T}}}
%\newcommand{\Xw}{\ensuremath{X_t(\omega)}}
%\newcommand{\Yw}{\ensuremath{Y_t(\omega)}}
%\newcommand{\Xq}{\ensuremath{X_q(\omega)}}
%\newcommand{\Yq}{\ensuremath{Y_q(\omega)}}
\newcommand{\E}{\ensuremath{\mathbb{E}}}
\newcommand{\1}{\ensuremath{\textbf{1}}}
\newcommand{\LwieLp}{\ensuremath{\mathcal{L}}}
% \L gibts anscheinend schon
\newcommand{\scriptL}{\ensuremath{\mathcal{L}}}
\newcommand{\scriptP}{\ensuremath{\mathcal{P}}}
\newcommand*\diff{\mathop{}\!\mathrm{d}}
\newcommand{\F}{\ensuremath{\mathcal{F}}}
\newcommand{\Dn}{\ensuremath{D_n}}
\newcommand{\cn}{\ensuremath{c_n}}
\newcommand{\ep}{\ensuremath{{e^{tc_n}}}}
\newcommand{\en}{\ensuremath{{e^{-tc_n}}}}
\newcommand{\N}{\ensuremath{\mathbb{N}}}
%\newcommand{\Xn}{\ensuremath{(X_n)_{n\in \N}}}
%\newcommand{\Yn}{\ensuremath{(Y_n)_{n\in \N}}}
\newcommand{\var}{\text{var}}
\newcommand{\textP}{\text{P}}


\begin{document}
	
\subsection*{Aufgabe 12}
\subsubsection*{Beispiel 2.2.8}

Seien $Y=(Y_i)_{i\ge 1}$ unabhängige, $\R$-wertige Zufallsvariablen auf $(\Om,\A,P)$ mit $Y_i\in\scriptL_{2}(P)$ und $\E Y_i=0$ für alle $i\ge 1$. Wir setzen $\F_0:=\{\emptyset,\Om\}$ und $\F_n:=\sigma(Y_1,\dots,Y_n)$ für $n\ge 1$. Dann definiert
$$X_n:=\sum_{i=1}^n Y_i,\qquad n\ge 0$$
nach \textit{Beispiel 2.1.7} ein Martingal bezüglich der Filtration $\F:=(\F_n)_{n\ge 0}$.
Ferner gilt $X_n\in \scriptL_2(P)$ für alle $n\ge 0$ und der quadratische Variationsprozess ist durch
$$\left<X\right>_n=\sum_{i=1}^n \E Y_i^2,\qquad n\ge 0$$
gegeben.

\subsubsection*{Beweis}
\einruecken

Da $\scriptL_2(P)$ ein Vektorraum ist, gilt $X_n=\sum_{i=1}^n Y_i\in\scriptL_2(P)$ für alle $n\ge 0$. Nach MWT-Vorlesung ist wegen $Y_i$ auch $Y_i^2$ unabhängig von $\F_{i-1}=\sigma(Y_1,\dots,Y_i-1)$, sodass
$$\E(Y_i|\F_{i-1})=\E(Y_i)$$
gilt. Nun erhalten wir aus \textit{Korollar 2.2.7 iv)} für $n\ge 0$:
\begin{align*}
\left<X\right>_n
&= \sum_{i=1}^n\E((X_i-X_{i-1})^2|\F_{i-1}) \\
&= \sum_{i=1}^n\E(Y_i^2|\F_{i-1}) \\
&= \sum_{i=1}^n\E Y_i^2.
\end{align*}

\ausruecken

\subsubsection*{Beispiel 2.2.9}

Seien $Y=(Y_i)_{i\ge 1}$ unabhängige, $\R$-wertige Zufallsvariablen auf $(\Om,\A,P)$ mit $Y_i \ge 0$, $Y_i\in\scriptL_{2}(P)$ und $\E Y_i=1$ für alle $i\ge 1$. Wir definieren $X_0:=1$, $\F_0:=\{\emptyset,\Om\}$ und
\begin{align*}
X_n &:= Y_1\cdot\dots\cdot Y_n \\
\F_n &:= \sigma(Y_1,\dots,Y_n)
\end{align*}
für $n\ge 1$. Dann ist $X:=(X_n)_{n\ge 0}$ nach \textit{Beispiel 2.1.9} ein Martingal bezüglich der Filtration $\F=(\F_n)_{n\ge 0}$ und es gilt $X_n\in\scriptL_2(P)$ für alle $n\ge0$. Ferner ist der quadratische Variationsprozess durch
$$\left<X\right>_n=\sum_{i=1}^n X_{i-1}^2\var Y_i,\qquad n\ge 0$$
gegeben.

\subsubsection*{Beweis}
\einruecken

Da $(Y_i)_{i\ge 1}$ unabhängig, sind auch $(|Y_i|^2)_{i\ge 1}$ unabhängig mit $|Y_i|^2\in\scriptL_1$.
Nach \textit{Satz 5.4.} aus MWT ist somit $|X_n|^2=\prod_{i=1}^n |Y_i|^2$ auch P-integrierbar, also $X_n\in\scriptL_2$.
Ferner gilt wie in letztem Beispiel auch $(Y_i-1)^2_{i\ge 1}$ unabhängig von $\F_{i-1}$, sodass
$$\E((Y_i-1)^2|\F_{i-1})=\E(Y_i-1)^2=\E(Y_i-\E Y_i)^2=\var Y_i$$

Nach \textit{Korollar 2.2.7 iv)} erhalten wir somit für $n\ge 0$:
\begin{align*}
\left<X\right>_n
&= \sum_{i=1}^n \E((X_i-X_{i-1})^2|\F_{i-1}) \\%,\qquad n\ge 0 \\
&= \sum_{i=1}^n \E((X_{i-1}Y_i - X_{i-1})^2|\F_{i-1}) \\
&= \sum_{i=1}^n \E(X_{i-1}^2\cdot(Y_i-1)^2|\F_{i-1}) \\
&= \sum_{i=1}^n X_{i-1}^2\cdot\E((Y_i-1)^2|\F_{i-1}) \tag{$\ast$} \\
&= \sum_{i=1}^n X_{i-1}^2\var Y_i,
\end{align*}
wobei $(\ast)$ wegen der $\F$-Adaptiertheit von $X$ gilt.

\ausruecken

\subsection*{Aufgabe 13}

Sei \WraumAnn, $\F=(\F_n){n\ge 0}$ eine Filtration, $(X_n)_{n\ge 0}$ ein $\F$-adaptierter Prozess und $a<b$. Definiere $\sigma_0:=0$ und, für $k\ge 1$,
\begin{align*}
\tau_k &:= \inf\{n\ge\sigma_{k-1}:X_n\le a\} \\
\sigma_k &:= \inf\{n\ge\tau_k:X_n\ge b\}.
\end{align*}
Zeige, dass $\sigma_{k-1}$ und $\tau_k$ $\F$-Stoppzeiten sind, die $\tau_k\le\sigma_k\le\tau_{k+1}$ für alle $k\ge 1$ erfüllen.
Wann sind diese Ungleichungen strikt?

\subsubsection*{Beweis und Lösung}
\einruecken

trivial

\ausruecken

\subsection*{Aufgabe 14}

Sei \WraumAnn, $\F=(\F_n)_{n\ge 0}$ eine Filtration und $B_n\in\F_n$ für alle $n\ge0$. Bestimme die quadratische Variation von dem Prozess gegeben durch
$$X_n:=\sum_{i=1}^n \1_{B_i},\qquad n\ge 0.$$

\subsubsection*{Lösung}
\einruecken

trivial

\ausruecken

\subsection*{Aufgabe 15}

Sei $X:=\Xt$ ein \R-wertiger stochastischer Prozess auf $(\Om,\A,P)$, sodass $X_t\in\scriptL_2(P)$ für alle $t\in T$. Wir definieren die Kovarianz-Funktion $k:T\times T \rightarrow \R$ durch
$$k(s,t):=\E((X_s-\E X_s)(X_t-\E X_t)),\qquad s,t\in T.$$
Zeige, dass die quadratische Matrix $(k(t_i,t_j))_{i,j=1}^n$ symmetrisch und positiv semi-definit ist für alle $n\ge 1$ und $t_1,\dots,t_n\in T$. Was können wir über die entsprechenden Eigenwerte sagen?

\subsubsection*{Beweis}
\einruecken

trivial

\ausruecken
	
\end{document}